\clearpage\SourceLeaf{184}
\begin{center}\textsc{Gauss} an \textsc{Olbers}. Göttingen, 28. April 1817.\end{center}

.\ .\ .\ .\ .\ .\ .\ \textsc{Wachter} hat eine kleine Schrift über die
ersten Grundlagen der Geometrie drucken lassen. Durch \textsc{Lindenau} werden Sie vermutlich
ein Exemplar erhalten. Obwohl W. tiefer in das Wesen der Sache
eingedrungen ist als seine Vorgänger, ist sein Beweis doch nicht
schlüssiger als alle anderen. Ich komme immer mehr zu der Überzeugung, dass die
Notwendigkeit unserer Geometrie nicht bewiesen werden kann, wenigstens
nicht \textls{vom} \textbf{\textls{menschlichen}} Verstand noch \textls{für} den menschlichen Verstand.
Vielleicht gelangen wir in einem anderen Leben zu anderen Einsichten in das
Wesen des Raums, die uns jetzt unerreichbar sind. Bis dahin müsste man die
Geometrie nicht der Arithmetik gleichstellen, die auf rein apriorischen Grundlagen beruht, sondern sie etwa
auf dieselbe Stufe wie die Mechanik stellen. .\ .\ .\ .\ .\ .

\begin{center}\rule{3cm}{0.4pt}\end{center}

